%%%%%%%%%%%%%%%%%%%%%%%%%%%%% Define Article %%%%%%%%%%%%%%%%%%%%%%%%%%%%%%%%%%
\documentclass{article}
%%%%%%%%%%%%%%%%%%%%%%%%%%%%%%%%%%%%%%%%%%%%%%%%%%%%%%%%%%%%%%%%%%%%%%%%%%%%%%%

%%%%%%%%%%%%%%%%%%%%%%%%%%%%% Using Packages %%%%%%%%%%%%%%%%%%%%%%%%%%%%%%%%%%
\usepackage{listings, xcolor}
\usepackage{parskip}
\usepackage{hyperref}
\usepackage{amsmath}
\usepackage{graphicx}
\usepackage[a4paper, margin = 1.5cm]{geometry}
\usepackage{float}
\usepackage{amssymb}

%%%%%%%%%%%%%%%%%%%%%%%%%%%%%%% Title & Author %%%%%%%%%%%%%%%%%%%%%%%%%%%%%%%%
\title{Social Computing}
\author{Trentin Tommaso}
%%%%%%%%%%%%%%%%%%%%%%%%%%%%%%%%%%%%%%%%%%%%%%%%%%%%%%%%%%%%%%%%%%%%%%%%%%%%%%%

\begin{document}
  \section{Parte 1}
    \subsection{Social Media}
      Gruppo di applicazioni internet-based che si basano sulle fondamenta ideologico-tecnologiche di Web 2.0, si definiscono \textbf{siti social network} i servizi web-based che permettono a individui di:
      \begin{itemize}
        \item costruire profili pubblici o semi-pubblici in un sistema bounded; 
        \item articolare una lista di altri utenti con cui condividere una connessione;
        \item vedere loro liste di connessioni e quelle fatte da altri utenti nel sistema.
      \end{itemize}
      \subsubsection{Classificazione}
        \begin{itemize}
          \item \textbf{Social presence / Media richness}: quanto "sembra che uno sia effettivamente lì"; 
          \item \textbf{Self presentation / Self disclosure}: quanto si rivela di sé stessi.
        \end{itemize}
        *Inserisci immagine*
      \subsubsection{Modello ad Alveare}
        *Inserisci Immagine*
      \subsubsection{Storia}
        
\end{document}
